%% kapa: exam=1 round=20 year=2009 period=2 form=A subject=appraisal_law subject_ko=부동산관계법규 seq=36 number=76
%% source: Q-Net 감정평가사 기출문제 (cst00309) artlSeq=5024178 file=extracted/1008/5024178/2009년 제20회 감정평가사 1차시험 기출문제/2교시.hwp
%% conversion: kordoc 4.24.51 -> tools/gen1.py (auto); review=proofread-sonnet; flags=none
\begin{question}{76}
지적법령상 지목에 관한 설명으로 옳지 \uline{않은} 것은?
\begin{choices}
  \item 지목은 법정지목인 철도용지를 포함하여 모두 28개 종류로 구분한다.
  \item 지적도에 등록된 사항 중 “원”으로 표기된 부호는 “과수원”이라는 법정지목을 의미하는 것이다.
  \item 토지가 일시적 또는 임시적인 용도로 사용되는 때에는 지목변경을 하지 않는다.
  \item 소관청은 토지의 이동현황을 직권조사하여 지목을 결정할 때 토지이동현황조사계획을 수립하여야 한다.
  \item 임야도에 등록된 토지가 사실상 형질변경 되었으나 지목변경을 할 수 없는 경우에는 지목의 변경 없이 등록전환을 신청할 수 있다.
\end{choices}
\end{question}
\answer{2}
